%% Gradient-controlled level-set transport on unstructured finite-volume meshes:
%% halo-limited velocity extension and admissible-band source terms.
%% Build:  latexmk -pdf gclsLevelSet.tex      (or: make article-gcls)
%%
%% STATUS: skeleton (2026-09-26). The continuum model (sections 3 and 4) follows
%% the two technical dossiers cited in the README of this theme. Sections 6-8 stay
%% empty until the gates named in them report; no number enters this file unless a
%% committed script wrote it into data/tables or data/figures.
%% Author list: to be decided by the authors.
\documentclass[preprint,3p,11pt]{elsarticle}

\usepackage{amsmath,amssymb}
\usepackage{graphicx}
\usepackage{booktabs}
\usepackage{url}

\graphicspath{{data/figures/}}

\newcommand{\psilev}{\psi}
\newcommand{\nhat}{\mathbf{n}}
\newcommand{\bu}{\mathbf{u}}
\newcommand{\bx}{\mathbf{x}}
\newcommand{\Dt}{\mathrm{D}_t}

\begin{document}
\begin{frontmatter}

\title{Gradient-controlled level-set transport on unstructured finite-volume
meshes: halo-limited velocity extension and admissible-band source terms}

\author[tuda]{Tomislav Mari\'{c}}
\author[tuda]{Dieter Bothe}
\address[tuda]{Mathematical Modeling and Analysis, TU Darmstadt, Germany}

\begin{abstract}
%% To be written after the 2D and 3D method gates report.
A signed-distance level set gives a well-scaled normal and a fixed physical
width of the interface transition, but the level-set transport equation does
not keep the signed-distance property under a deforming flow. This work keeps
the gradient magnitude $q=|\nabla\psilev|$ inside an admissible band near the
interface instead of enforcing $q=1$, with no redistancing, on unstructured
finite-volume meshes under MPI domain decomposition.
\end{abstract}

\begin{keyword}
level set \sep velocity extension \sep signed distance \sep unstructured finite
volume \sep OpenFOAM \sep two-phase flow
\end{keyword}

\end{frontmatter}

%% ---------------------------------------------------------------------------
\section{Introduction}
\label{sec:intro}
%% The problem: D_t q = -a q; why redistancing is unwanted (interface
%% displacement, cost, contact lines); the admissible band as the target.
For the material equation $\partial_t\psilev+\bu\cdot\nabla\psilev=0$ the
gradient magnitude obeys, along a trajectory,
\begin{equation}
  \Dt q = -a\,q,\qquad a=\nhat^{\mathsf T}(\nabla\bu)\,\nhat,\qquad
  \nhat=\nabla\psilev/q ,
  \label{eq:qlaw}
\end{equation}
so the normal strain $a$ changes $q$ even for an exact solution
\cite{FrickeEtAl2022,BotheFrickeSoga2024}.

%% ---------------------------------------------------------------------------
\section{Prior art}
\label{sec:priorart}
%% From the halo-limited dossier, section "State of the art": hyperbolic
%% reinitialisation \cite{SussmanSmerekaOsher1994,SussmanFatemi1999}; geometric
%% and Eikonal methods \cite{KarakusEtAl2022,Shakoor2025}; conservative level
%% set and CLSVOF \cite{OlssonKreiss2005,OlssonKreissZahedi2007,KeesEtAl2011,LyrasEtAl2020};
%% variational methods \cite{LiEtAl2010,ToureSoulaimani2016,ShaoEtAl2023};
%% modified equations \cite{SabelnikovEtAl2014,Hamamuki2019}; SDPLS and the
%% bounded-gradient equation \cite{FrickeEtAl2022,BotheFrickeSoga2024};
%% gradient-augmented level set \cite{NaveRosalesSeibold2010,BockmannVartdal2014};
%% velocity extension \cite{AdalsteinssonSethian1999,BotheSoga2026}; OpenFOAM
%% and LEIA \cite{WellerEtAl1998,LEIA,Reitzel2023,LiuEtAl2023,FerroEtAl2025}.

%% ---------------------------------------------------------------------------
\section{Continuum model}
\label{sec:model}

\subsection{Source calculus}
For $\Dt\psilev=\psilev F$ the zero level set is invariant, and on it
\begin{equation}
  \Dt q = q\,(F-a).
  \label{eq:designlaw}
\end{equation}

\subsection{Halo-limited directional velocity extension}
Let $d$ be a signed nominal distance and $\mathbf e$ the unit direction from
the interface toward $\bx$; the level-set realisation is $d=\psilev/q$,
$\mathbf e=\nhat$. With a radius $R$ and $m\in\mathbb N$,
\begin{equation}
  S_R(d)=d\left[1+\left(\frac{d}{R}\right)^{2m}\right]^{-1/(2m)},\qquad
  Y_R=\bx-S_R(d)\,\mathbf e,\qquad
  w_R=\left(\frac{S_R(d)}{d}\right)^{\beta},
\end{equation}
\begin{equation}
  \bu_H=(1-w_R)\,\bu(\bx)+w_R\,\bu(Y_R).
\end{equation}
Every sample lies inside the radius $R$, and on the interface
$\nhat^{\mathsf T}(\nabla\bu_H)\nhat=0$.

\subsection{Source-law family}
%% Table: none, linearQ mu(1-q), linearZ -mu (q^2-1)/2, cubic laws,
%% regularised 2/3 law, saturated linear z law, the soft wall
%% -kappa tanh[gamma((q-1)/delta_s)^p]; the strain weight none|full|omega;
%% the robust inward-pointing condition kappa Theta* >= sigma + ... + mu.

%% ---------------------------------------------------------------------------
\section{Discretisation on unstructured finite-volume meshes}
\label{sec:fv}
%% The material-form operator div_h(Phi psi) - psi div_h(Phi) with the ACTUAL
%% transport flux; the two discrete divergences; the flux correction
%% Phi^H = Phi^NS + w [U_h(Y) - U_h(x)].S_f; the partition-independent stencil
%% sampler; the exponential source update; the semi-Lagrangian realisation
%% (trace along reconstruct(Phi^H), split source step).

%% ---------------------------------------------------------------------------
\section{Implementation in leia}
\label{sec:impl}
%% Runtime-selectable families (velocityExtension, sdplsSource, slSource,
%% gradientControlLaw and the strategies of haloLimited), composition roots per
%% solver, extension without modification, the parallel design.

%% ---------------------------------------------------------------------------
\section{Verification}
\label{sec:verification}
%% EMPTY until the gates report: unit and exact gates (the law closed forms,
%% the interfacial cancellation, constant-state preservation), the static
%% gradient metrology E0, the exact 1D gates E1, the seam gate E2.

%% ---------------------------------------------------------------------------
\section{Coupled results}
\label{sec:coupled}
%% EMPTY until the method gates report: methodGate2D and methodGate3D tables
%% (data/tables/<gate>_<candidate>_{summary,orders,vsBaseline}.csv) with the
%% observed orders and the GCI \cite{Roache1994,Celik2008}.

%% ---------------------------------------------------------------------------
\section{Parallel cost and decomposition invariance}
\label{sec:parallel}
%% EMPTY until the gates report: rank count, wall clock, core-seconds per
%% step, serial vs np 4 vs np 8 differences.

%% ---------------------------------------------------------------------------
\section{Discussion}
\label{sec:discussion}
%% What was falsified; limits (volume conservation, transition strain).

%% ---------------------------------------------------------------------------
\section{Conclusions}
\label{sec:conclusions}

\section*{Acknowledgements}
%% To be written.

\bibliographystyle{elsarticle-num}
\bibliography{refs}
\end{document}
